\section{Introduction}

Multi-agent architectures built from large language models are now routine
engineering practice. Frameworks such as AutoGen, MetaGPT and ChatDev
\citep{wu2023autogen,hong2024metagpt,qian2024chatdev} let a practitioner
decompose a task across several model instances that propose, critique and
aggregate one another's work, and reported gains over a single model are
common. Yet the picture from the literature is contradictory. Multi-agent
debate improves factuality and reasoning on some benchmarks
\citep{du2023debate}; simply sampling more answers and voting captures much of
the same benefit \citep{wang2023selfconsistency,li2024moreagents}; and a recent
study of deployed multi-agent systems finds failure rates high enough to
question whether the additional structure is doing useful work at all
\citep{cemri2025why}. Practitioners deciding whether to build a multi-agent
system therefore have no principled way to anticipate the answer, and the
prevailing method---build it, then measure---is expensive and often
inconclusive.

We argue that two problems explain most of this disagreement, and that both
are fixable.

\textbf{The baseline is usually not compute-matched.} A system of $k$ agents
spends roughly $k$ times the tokens of a single agent. Comparing it against one
agent that produced one answer credits the architecture for compute that any
method could have spent. The honest comparison prices the single agent at the
same budget---letting it think longer, or revise its own answer---and asks
whether \emph{coordination}, rather than \emph{compute}, produced the gain.
This concern has been raised for agent evaluation generally
\citep{kapoor2024agents} and for repeated sampling in particular
\citep{brown2024monkeys,snell2024scaling,stroebl2024limits}, but
compute-matched comparison is still not standard practice in the multi-agent
literature.

\textbf{There is no measurable quantity that says when to expect a gain.}
Existing accounts are largely taxonomic: they describe what multi-agent systems
do, or catalogue how they fail, without predicting where gains will appear.
What is missing is a property of the \emph{task and model together} that can be
measured cheaply, in advance, and that determines the answer.

This paper supplies that quantity. Our starting point is an observation that is
almost trivial once stated but that organises the entire empirical picture: an
ensemble of $k$ agents can only ever be as good as the best candidate it
produces, and it realises that potential only to the extent that its selection
mechanism can tell good candidates from bad. Writing $\pone$ for single-sample
accuracy, $\cov(k)$ for the probability that at least one of $k$ samples is
correct, and $A_\sigma(k)$ for the accuracy of a system using selector $\sigma$,
\begin{equation}
A_\sigma(k) \;=\; \pone + \eff_\sigma(k)\,\bigl(\cov(k) - \pone\bigr),
\label{eq:decomp-intro}
\end{equation}
where $\eff_\sigma(k) \in [0,1]$ is the fraction of the available
\emph{coverage headroom} $\cov(k)-\pone$ that the selector converts into
accuracy. Multi-agent gains require both factors to be non-trivial: headroom
without a discriminating selector is wasted, and a perfect selector on a pool
with no headroom has nothing to find.

The useful consequence is that both factors can be estimated \emph{before}
building the system, from a small labelled calibration set. Coverage headroom
needs only repeated sampling. For selection we introduce the
\textbf{generator--verifier gap} $\gap$: the probability, in excess of chance,
that the model prefers a correct candidate over an incorrect one when shown
both. Because $\gap$ is a \emph{pairwise} quantity while selection is
\emph{listwise}, we supply the link with a latent-score model
(Section~\ref{sec:framework}) that turns a measured $\gap$ into a
parameter-free prediction of $\eff$. Verifier-free aggregation such as
self-consistency needs no verifier at all and is governed by a second
diagnostic, \textbf{plurality alignment} $\plur$, which measures how well the
modal answer tracks the truth.

We test this on \nTasks{} tasks spanning six families chosen to span the
verification spectrum---from constraint-satisfaction problems whose answers can
be checked in a line of code, to transitive-ordering puzzles where confirming
an answer means redoing the derivation---across \nModels{} model scales and
\nArch{} architectures, all with exact token accounting. Our contributions:

\begin{enumerate}[leftmargin=1.4em,itemsep=2pt,topsep=3pt]
\item \textbf{A decomposition and a theory.} Equation~\eqref{eq:decomp-intro}
  separates a multi-agent system's accuracy into coverage and selection
  efficiency, and a latent-score model predicts selection efficiency from the
  pairwise gap $\gap$ with no free parameters (Section~\ref{sec:framework}).
\item \textbf{Token-matched evidence that most multi-agent gains are compute, not
  coordination.} Against a single agent priced at the same token budget,
  \nFlip{} of the architecture--task cells that appear to gain lose that gain,
  and the median configuration moves from \medianNaiveLift{} to
  \medianMatchedLift{}---a change of sign (Section~\ref{sec:results}).
\item \textbf{A prospective diagnostic.} Measured on \nCalib{} calibration
  tasks and passed through the model with no fitted parameters, the
  diagnostics predict held-out lift with $R^2 = \rTwoApriori{}$ without the
  multi-agent system ever being built. As a binary build/don't-build rule they
  are weaker---\decisionAgreement{} correct against a \decisionBaseline{}
  majority baseline---and we report that gap rather than the raw accuracy.
\item \textbf{A case study in a deployed setting.} On citation generation,
  where fabricated output is a known failure mode, an external registry check
  converts nearly all available headroom at zero token cost while the same
  model used as its own critic performs near chance
  (Section~\ref{sec:casestudy}).
\end{enumerate}

Our headline practical finding is a negative result with a constructive
remedy. Multi-agent structure is not a general-purpose accuracy technique; it
is a way of spending compute that pays off precisely when an asymmetry exists
between producing an answer and checking one. Where that asymmetry is absent,
the same tokens are better spent on a single agent---and where it is present,
the effective move is usually not another language-model critic but an external
check, which in our experiments costs no tokens at all.
